% ════════════════════════════════════════════
% 局限与改进方向（08.模型改进推广.tex）
% ════════════════════════════════════════════
\section{局限与改进方向}

\subsection{研究局限与改进}

当前模型的主要不足来自数据维度而非估计方法。若条件允许，可补充帖子级的时间戳、受众画像、互动数据与内容文本，把同一主页的帖子组织成面板，用帖子固定效应吸收不随时间变化的内容差异，再用月份固定效应控制共同冲击，从而把付费效应从账号成长与季节波动中更干净地分离出来。对视频这类稀疏类别，可采用分层贝叶斯或收缩估计，把小组均值向总体均值适度拉近，以降低区间宽度；对曝光量的重尾，可比较 Tweedie 分布与分位数回归，前者直接匹配“零膨胀加连续正值”的计数特征，后者能回答不同曝光水平上的差异化效应。若要把相关关系推进到因果结论，还可利用付费投放的时间断点或倾向得分匹配构造准实验，检验本报告的平均效应与主题异质性是否稳定。

此外，截面数据无法控制帖子层面的固定差异，模型也未纳入推荐算法、受众画像与内容质量等关键变量，因此样本外拟合优度仅约 0.10。后续可在数据许可的前提下补充这些变量，或采用面板固定效应与工具变量策略，把“平均关联”进一步推进为“条件因果效应”。在样本量有限时，还应优先扩大视频与链接类别的观测，避免区间过宽导致结论只能停留在方向判断；主题异质性目前只有主题 1 与主题 3 的对照较为清晰，主题 2 处于边缘显著，也需要更多样本确认。

\subsection{方法推广}

本报告的“共线诊断—固定效应替代—烟雾还原—参数模拟”流程不依赖具体平台，可迁移到微博、短视频与电商直播等场景的曝光或流量分析，也可用于广告投放的增量评估与内容排期优化。在工程实践中，该流程可以作为营销数据看板的后端模型，按周更新系数并输出各内容类型的调整预测；在行业层面，它提供了一种在账号规模与时间趋势高度绑定的条件下仍能给出可解释效应的分析范式，对预算分配与内容策略具有参考价值。
